% GENERATED FILE. Edit canonical lessons, then run npm run generate:books.
% canonical-source-hash: fnv1a64:48366eeb6160cba4
% canonical-lessons: GE-C529-renovieren, GE-C529-tapete, GE-C529-tapezieren, GE-C529-fliese, GE-C529-dachboden, GE-R529-review-522-529

\chapter{Renovieren}
\label{ch:ge-renovieren}
\hlchaptermodality{529}

\begin{chapteropening}
\textbf{By the end of this chapter:} \emph{I can talk about doing up a flat.}

Complete the last lesson of chapter 529: I can talk about doing up a flat.
\end{chapteropening}

\section[renovieren]{renovieren — to renovate, to do up}

\label{lesson:GE-C529-renovieren}



\begin{quote}
Before the new one: say the German for dustbin, then the German for to dispose of.
\end{quote}

\subsection*{You'll want to know: renovieren}
\textbf{renovieren} — \textquotedblleft{}to renovate, to do up\textquotedblright{}.

\begin{quote}
Vor dem Einzug haben wir die Wohnung renoviert.
\end{quote}

\begin{quote}
\emph{Before moving in, we did up the flat.}
\end{quote}

Regular, Perfekt without \emph{ge-}: \emph{hat renoviert}.

\begin{cousinweb}
\emph{renovieren} was borrowed from Middle French \emph{renover} (EtymDB), from Latin \emph{renovāre} \textquotedblleft{}to make new\textquotedblright{}, built on \emph{novus} \textquotedblleft{}new\textquotedblright{}.
\end{cousinweb}

\begin{grammarlens}[title={What you've built}]
One more word for telling, explaining and describing at B1.
\end{grammarlens}

\subsection*{Your turn}
\begin{itemize}
\raggedright
  \item \emph{Say it:} \emph{renovieren}
  \item \emph{Say it:} \emph{renovieren}, once more
  \item \emph{Say it:} a sentence of your own with \emph{renovieren}
  \item \emph{Recall:} say the German for dustbin, then the German for to dispose of, then say \emph{renovieren} again
\end{itemize}

\begin{tcolorbox}[breakable,colback=teal!4,colframe=teal!35!black,title={Before you move on}]
What does renovieren mean? (\textquotedblleft{}To renovate, to do up\textquotedblright{}.) Say it once more.
\end{tcolorbox}
\section[die Tapete]{die Tapete — wallpaper}

\label{lesson:GE-C529-tapete}



\begin{quote}
Before the new one: say the German for to dispose of, then the German for to renovate.
\end{quote}

\subsection*{You'll want to know: die Tapete}
\textbf{die Tapete} — \textquotedblleft{}wallpaper\textquotedblright{}.

\begin{quote}
Die alte Tapete im Flur ist hässlich.
\end{quote}

\begin{quote}
\emph{The old wallpaper in the hall is ugly.}
\end{quote}

Plural: \emph{die Tapeten}.

\begin{cousinweb}
\emph{Tapete} comes from Medieval Latin \emph{tapeta} \textquotedblleft{}wall carpet\textquotedblright{} (EtymDB), from Latin \emph{tapēte} \textquotedblleft{}carpet\textquotedblright{}.
\end{cousinweb}

\begin{grammarlens}[title={What you've built}]
One more word for telling, explaining and describing at B1.
\end{grammarlens}

\subsection*{Your turn}
\begin{itemize}
\raggedright
  \item \emph{Say it:} \emph{die Tapete}
  \item \emph{Say it:} \emph{die Tapete}, once more
  \item \emph{Say it:} a sentence of your own with \emph{die Tapete}
  \item \emph{Recall:} say the German for to dispose of, then the German for to renovate, then say \emph{die Tapete} again
\end{itemize}

\begin{tcolorbox}[breakable,colback=teal!4,colframe=teal!35!black,title={Before you move on}]
What does die Tapete mean? (\textquotedblleft{}Wallpaper\textquotedblright{}.) Say it once more.
\end{tcolorbox}
\section[tapezieren]{tapezieren — to wallpaper}

\label{lesson:GE-C529-tapezieren}



\begin{quote}
Before the new one: say the German for to renovate, then the German for wallpaper.
\end{quote}

\subsection*{You'll want to know: tapezieren}
\textbf{tapezieren} — \textquotedblleft{}to wallpaper\textquotedblright{}.

\begin{quote}
Am Wochenende tapeziert meine Schwester das Kinderzimmer.
\end{quote}

\begin{quote}
\emph{At the weekend my sister is wallpapering the children's room.}
\end{quote}

Regular, Perfekt without \emph{ge-}: \emph{hat tapeziert}.

\begin{cousinweb}
\emph{tapezieren} is built on \emph{Tapete} \textquotedblleft{}wallpaper\textquotedblright{}, from Medieval Latin \emph{tapeta} \textquotedblleft{}wall carpet\textquotedblright{} (EtymDB).
\end{cousinweb}

\begin{grammarlens}[title={What you've built}]
One more word for telling, explaining and describing at B1.
\end{grammarlens}

\subsection*{Your turn}
\begin{itemize}
\raggedright
  \item \emph{Say it:} \emph{tapezieren}
  \item \emph{Say it:} \emph{tapezieren}, once more
  \item \emph{Say it:} a sentence of your own with \emph{tapezieren}
  \item \emph{Recall:} say the German for to renovate, then the German for wallpaper, then say \emph{tapezieren} again
\end{itemize}

\begin{tcolorbox}[breakable,colback=teal!4,colframe=teal!35!black,title={Before you move on}]
What does tapezieren mean? (\textquotedblleft{}To wallpaper\textquotedblright{}.) Say it once more.
\end{tcolorbox}
\section[die Fliese]{die Fliese — tile}

\label{lesson:GE-C529-fliese}



\begin{quote}
Before the new one: say the German for wallpaper, then the German for to wallpaper.
\end{quote}

\subsection*{You'll want to know: die Fliese}
\textbf{die Fliese} — \textquotedblleft{}tile\textquotedblright{}.

\begin{quote}
Im Bad sind die Fliesen weiß.
\end{quote}

\begin{quote}
\emph{In the bathroom the tiles are white.}
\end{quote}

Plural: \emph{die Fliesen}. Also \emph{die Kachel}.

\begin{cousinweb}
\emph{Fliese} was borrowed from Middle Low German \emph{vlīse} (EtymDB), a word of the north.
\end{cousinweb}

\begin{grammarlens}[title={What you've built}]
One more word for telling, explaining and describing at B1.
\end{grammarlens}

\subsection*{Your turn}
\begin{itemize}
\raggedright
  \item \emph{Say it:} \emph{die Fliese}
  \item \emph{Say it:} \emph{die Fliese}, once more
  \item \emph{Say it:} a sentence of your own with \emph{die Fliese}
  \item \emph{Recall:} say the German for wallpaper, then the German for to wallpaper, then say \emph{die Fliese} again
\end{itemize}

\begin{tcolorbox}[breakable,colback=teal!4,colframe=teal!35!black,title={Before you move on}]
What does die Fliese mean? (\textquotedblleft{}Tile\textquotedblright{}.) Say it once more.
\end{tcolorbox}
\section[der Dachboden]{der Dachboden — attic, loft}

\label{lesson:GE-C529-dachboden}



\begin{quote}
Before the new one: say the German for to wallpaper, then the German for tile.
\end{quote}

\subsection*{You'll want to know: der Dachboden}
\textbf{der Dachboden} — \textquotedblleft{}attic, loft\textquotedblright{}.

\begin{quote}
Die alten Spielsachen liegen auf dem Dachboden.
\end{quote}

\begin{quote}
\emph{The old toys are in the attic.}
\end{quote}

Plural: \emph{die Dachböden}.

\begin{cousinweb}
\emph{Dachboden} is built on \emph{Dach} \textquotedblleft{}roof\textquotedblright{}, inherited from Old High German \emph{dah}, from Proto-Germanic \emph{þaką} (EtymDB), and \emph{Boden} \textquotedblleft{}floor\textquotedblright{}.
\end{cousinweb}

\begin{grammarlens}[title={What you've built}]
One more word for telling, explaining and describing at B1.
\end{grammarlens}

\subsection*{Your turn}
\begin{itemize}
\raggedright
  \item \emph{Say it:} \emph{der Dachboden}
  \item \emph{Say it:} \emph{der Dachboden}, once more
  \item \emph{Say it:} a sentence of your own with \emph{der Dachboden}
  \item \emph{Recall:} say the German for to wallpaper, then the German for tile, then say \emph{der Dachboden} again
\end{itemize}

\begin{tcolorbox}[breakable,colback=teal!4,colframe=teal!35!black,title={Before you move on}]
What does der Dachboden mean? (\textquotedblleft{}Attic, loft\textquotedblright{}.) Say it once more.
\end{tcolorbox}
\section[(dialogue)]{Words from chapters 522–529 — a review}

\label{lesson:GE-R529-review-522-529}



\begin{quote}
Say the German for to borrow and to leaf.
\end{quote}

\subsection*{Your turn}
Say these aloud:

\begin{itemize}
\raggedright
  \item to borrow, to leaf, bookmark, publisher, translation
  \item remote control, to flick through channels, series, season, ratings
  \item comedy, tragedy, documentary, subtitle, to dub
  \item to stream, to subscribe to, video game, to play video games, headphones
  \item dust, to wipe, to sweep, to vacuum, cloth
  \item tumble dryer, clothes line, clothes peg, to iron, to fold
  \item mess, to clear out, rubbish, dustbin, to dispose of
  \item to renovate, wallpaper, to wallpaper, tile, attic
\end{itemize}

\begin{tcolorbox}[breakable,colback=teal!4,colframe=teal!35!black,title={Before you move on}]
To borrow? (\textbf{ausleihen}.) To leaf? (\textbf{blättern}.) Bookmark? (\textbf{das Lesezeichen}.) Publisher? (\textbf{der Verlag}.) Translation? (\textbf{die Übersetzung}.) Remote control? (\textbf{die Fernbedienung}.) To flick through channels? (\textbf{zappen}.) Series? (\textbf{die Serie}.) Season? (\textbf{die Staffel}.) Ratings? (\textbf{die Einschaltquote}.) Comedy? (\textbf{die Komödie}.) Tragedy? (\textbf{die Tragödie}.) Documentary? (\textbf{der Dokumentarfilm}.) Subtitle? (\textbf{der Untertitel}.) To dub? (\textbf{synchronisieren}.) To stream? (\textbf{streamen}.) To subscribe to? (\textbf{abonnieren}.) Video game? (\textbf{das Videospiel}.) To play video games? (\textbf{zocken}.) Headphones? (\textbf{der Kopfhörer}.) Dust? (\textbf{der Staub}.) To wipe? (\textbf{wischen}.) To sweep? (\textbf{fegen}.) To vacuum? (\textbf{staubsaugen}.) Cloth? (\textbf{der Lappen}.) Tumble dryer? (\textbf{der Trockner}.) Clothes line? (\textbf{die Wäscheleine}.) Clothes peg? (\textbf{die Wäscheklammer}.) To iron? (\textbf{bügeln}.) To fold? (\textbf{falten}.) Mess? (\textbf{die Unordnung}.) To clear out? (\textbf{ausmisten}.) Rubbish? (\textbf{der Abfall}.) Dustbin? (\textbf{die Mülltonne}.) To dispose of? (\textbf{entsorgen}.) To renovate? (\textbf{renovieren}.) Wallpaper? (\textbf{die Tapete}.) To wallpaper? (\textbf{tapezieren}.) Tile? (\textbf{die Fliese}.) Attic? (\textbf{der Dachboden}.)
\end{tcolorbox}
